\documentclass[12pt, a4paper]{article}

% --- PACOTES ---
\usepackage[utf8]{inputenc}
\usepackage[T1]{fontenc}
\usepackage[brazilian]{babel}
\usepackage[left=1.5cm, right=1.5cm, top=1.5cm, bottom=1.5cm, footskip=0.5cm]{geometry}
\usepackage{amsmath, amsfonts, amssymb}
\usepackage{graphicx}
\usepackage{xcolor}
\usepackage{tikz}
\usepackage{eso-pic}
\usepackage{listings} % Pacote para formatar os códigos Java

% ==========================================
%        CONFIGURAÇÃO PARA CÓDIGO JAVA
% ==========================================
\definecolor{codegreen}{rgb}{0,0.6,0}
\definecolor{codegray}{rgb}{0.5,0.5,0.5}
\definecolor{codepurple}{rgb}{0.58,0,0.82}
\definecolor{backcolour}{rgb}{0.95,0.95,0.92}

\lstdefinestyle{mystyle}{
    backgroundcolor=\color{backcolour},   
    commentstyle=\color{codegreen},
    keywordstyle=\color{magenta},
    numberstyle=\tiny\color{codegray},
    stringstyle=\color{codepurple},
    basicstyle=\ttfamily\footnotesize,
    breakatwhitespace=false,         
    breaklines=true,                 
    captionpos=b,                    
    keepspaces=true,                 
    numbers=left,                    
    numbersep=5pt,                  
    showspaces=false,                
    showstringspaces=false,
    showtabs=false,                  
    tabsize=4,
    language=Java
}
\lstset{style=mystyle}

% ==========================================
%        CAMPOS EDITÁVEIS (FORMULÁRIO)
% ==========================================

\newcommand{\instituicao}{INSTITUTO FEDERAL DE SERGIPE}
\newcommand{\campus}{CAMPUS LAGARTO}
\newcommand{\curso}{BACHARELADO EM SISTEMAS DE INFORMAÇÃO}

\newcommand{\professor}{Jose Solenir Lima Figueredo}
\newcommand{\disciplina}{Algoritmo e Estrutura de Dados I}
\newcommand{\monitor}{José Gustavo Correia Nascimento}

\newcommand{\data}{\underline{\hspace{0.7cm}}/\underline{\hspace{0.7cm}}/\underline{\hspace{1cm}}}
\newcommand{\estudante}{\hrulefill}

\newcommand{\titulo}{LISTA DE EXERCÍCIOS DE REVISÃO}

\newcommand{\logoEsquerda}{geral-logo-ifs.png}

% ==========================================
%            BORDA DA PÁGINA
% ==========================================

\AddToShipoutPictureBG{
\begin{tikzpicture}[remember picture, overlay]
\draw[line width=1pt]
([xshift=0.5cm, yshift=-0.5cm]current page.north west)
rectangle
([xshift=-0.5cm, yshift=0.5cm]current page.south east);
\end{tikzpicture}
}

\begin{document}

% ==========================================
%                CABEÇALHO
% ==========================================

% --- LOGO E TÍTULO INSTITUCIONAL ---
\noindent
\begin{minipage}[c]{0.2\textwidth}
    % Se não tiver a logo no momento, você pode comentar a linha abaixo
    \includegraphics[width=\linewidth]{\logoEsquerda}
\end{minipage}%
\begin{minipage}[c]{0.8\textwidth}
    \centering
    \textbf{\instituicao}
    \textbf{\campus}\\[0.1cm]
    \textbf{\curso}
\end{minipage}%
\begin{minipage}[c]{0.2\textwidth}
    % Espaço vazio para equilibrar a logo da esquerda
\end{minipage}

\vspace{0.7cm}

% --- INFORMAÇÕES DA DISCIPLINA ---
\noindent
\textbf{Professor(a):} \professor \hfill \textbf{Disciplina:} \disciplina \\[0.2cm]
\noindent
\textbf{Monitor(a):} \monitor \hfill \textbf{Data:} \data \\[0.3cm]
\noindent
\textbf{Estudante:} \estudante

\vspace{0.5cm}

% --- TÍTULO ---
\begin{center}
    \Large\textbf{\titulo}
\end{center}
\vspace{0.3cm}
% ==========================================
%            FIM DO CABEÇALHO
% ==========================================

% --- EXEMPLO DE COMO USAR AS QUESTÕES E O CÓDIGO ---

\section*{Questão 1: Auditoria de Lotes (Conceito)}
Em uma fábrica de peças automotivas, os inspetores de qualidade registram a quantidade de peças defeituosas encontradas diariamente. Esses registros são armazenados em um vetor (array) de inteiros. Para automatizar o relatório semanal, você deve implementar um sistema que calcule o total de peças defeituosas utilizando \textbf{recursividade linear}.

\textbf{Requisitos:}
\begin{itemize}
    \item Implemente um método recursivo chamado \texttt{somarDefeitos} que receba o array de registros e o tamanho do array (ou o índice atual) como parâmetros.
    \item O método deve retornar a soma total dos elementos do array.
    \item Identifique e implemente claramente o \textbf{caso base} e o \textbf{passo recursivo}.
    \item Crie uma classe \texttt{Main} para testar o seu método com o array \texttt{\{2, 5, 1, 0, 4\}}.
\end{itemize}

\textbf{Pergunta Reflexiva:} Explique, com suas palavras, o que acontece com a pilha de execução (\textit{Call Stack}) do Java a cada nova chamada recursiva do método \texttt{somarDefeitos}. Como as variáveis locais e o estado de cada chamada são mantidos em memória até que o caso base seja atingido?

\vspace{0.5cm}

\section*{Questão 2: Validador de Códigos de Segurança (Prática)}
Um sistema de cofre digital utiliza códigos de acesso simétricos para liberar a abertura das portas. Um código é considerado válido apenas se for um palíndromo perfeito (ou seja, lido da mesma forma da esquerda para a direita e da direita para a esquerda, ignorando espaços e diferenças entre maiúsculas e minúsculas).

\textbf{Requisitos:}
\begin{itemize}
    \item Desenvolva um método recursivo \texttt{isPalindromo} que receba uma \texttt{String} (o código) e os índices de início e fim da análise.
    \item O método deve retornar um valor \texttt{boolean} indicando se o código é válido.
    \item Realize o tratamento para ignorar espaços em branco e letras maiúsculas/minúsculas da \texttt{String} original na classe principal antes de passar a informação para o método recursivo.
    \item A recursão deve comparar os caracteres das extremidades e se aproximar do centro a cada chamada (recursão de cauda/linear).
    \item Teste sua implementação com os códigos \texttt{"Ovo"}, \texttt{"Ame o poema"} e \texttt{"Senha123"}.
\end{itemize}

\textbf{Pergunta Reflexiva:} Em implementações recursivas como esta, a definição correta do caso base é vital. O que ocorreria nos bastidores (em termos de consumo de recursos e comportamento da JVM) se o caso base nunca fosse alcançado? Justifique.

\vspace{0.5cm}

\section*{Questão 3: Rastreador de Transações Financeiras (Desafio)}
Um banco processa milhões de transações diárias. Para otimizar o sistema de auditoria e detecção de fraudes, os IDs das transações (números inteiros) são armazenados em um array estritamente ordenado de forma crescente. Sua missão é implementar um motor de busca de alto desempenho para localizar rapidamente o índice de uma transação suspeita, seguindo o paradigma de divisão e conquista.

\textbf{Requisitos:}
\begin{itemize}
    \item Implemente um método recursivo de \textbf{Busca Binária} chamado \texttt{buscarTransacao} que receba o array de transações, o ID desejado e os limites do sub-array (início e fim).
    \item Se a transação for encontrada, retorne o índice correspondente no array.
    \item Se a transação não existir, retorne \texttt{-1}.
    \item Demonstre o funcionamento na classe \texttt{Main} utilizando um array contendo pelo menos 10 IDs ordenados e realize buscas por um ID existente e por um inexistente.
\end{itemize}

\textbf{Pergunta Reflexiva:} Ao utilizar a estratégia de divisão e conquista da Busca Binária recursiva, o tamanho do problema cai pela metade a cada chamada. Como essa característica afeta a profundidade máxima da pilha de chamadas (\textit{Call Stack}) em comparação com uma busca recursiva simples de ponta a ponta (busca linear) para o mesmo array? Explique a vantagem de desempenho com base nisso.

% ==========================================
%     TEMA 2: ARRAY PRIMITIVO (TAMANHO FIXO)
% ==========================================
\newpage

\noindent\rule{\textwidth}{0.8pt}
\begin{center}
    \large\textbf{TEMA 2 — ARRAY PRIMITIVO: ALOCAÇÃO CONTÍGUA E TAMANHO FIXO}
\end{center}
\noindent\rule{\textwidth}{0.8pt}
\vspace{0.4cm}

\section*{Questão 4: Controle de Vagas de Estacionamento (Conceito)}
Uma empresa de mobilidade urbana opera um estacionamento inteligente com um número fixo de vagas numeradas (de 0 a $N-1$). O sistema mantém um \textbf{array primitivo de Strings} de tamanho fixo, onde cada posição armazena a placa do veículo estacionado ou \texttt{null} se a vaga estiver livre. Você deve implementar um painel de controle básico para esse estacionamento.

\textbf{Requisitos:}
\begin{itemize}
    \item Declare um array de \texttt{String} com capacidade fixa de \textbf{8 vagas}.
    \item Implemente um método \texttt{estacionar(String[] vagas, int posicao, String placa)} que registre um veículo em uma posição específica. Se a posição já estiver ocupada, o método deve exibir um aviso e \textbf{não} sobrescrever a placa existente.
    \item Implemente um método \texttt{liberar(String[] vagas, int posicao)} que remova o veículo da posição informada (atribuindo \texttt{null}).
    \item Implemente um método \texttt{exibirPainel(String[] vagas)} que percorra todo o array e imprima o estado de cada vaga (livre ou ocupada com a respectiva placa).
    \item Na classe \texttt{Main}, simule a seguinte sequência: estacione 4 veículos em posições à sua escolha, tente estacionar um 5º veículo em uma posição já ocupada, libere uma vaga e exiba o painel final.
\end{itemize}

\textbf{Pergunta Reflexiva:} O array de vagas foi criado com tamanho fixo de 8 posições. Suponha que o estacionamento precise expandir para 12 vagas. Explique, do ponto de vista da alocação de memória contígua, por que \textbf{não é possível} simplesmente ``aumentar'' o array existente. Qual seria a estratégia necessária para acomodar as novas vagas e qual o custo computacional envolvido nessa operação?

\vspace{0.5cm}

\section*{Questão 5: Gerenciador de Playlist Musical (Prática)}
Uma plataforma de streaming de música permite que o usuário crie playlists com um limite máximo de músicas. Internamente, a playlist é representada por um \textbf{array primitivo de Strings} de capacidade fixa e uma variável \texttt{tamanho} que controla quantas músicas foram efetivamente inseridas. Você deve construir um gerenciador completo que demonstre as operações de inserção e remoção e seus custos.

\textbf{Requisitos:}
\begin{itemize}
    \item Crie uma classe \texttt{Playlist} que encapsule:
    \begin{itemize}
        \item Um array \texttt{String[] musicas} com capacidade fixa de \textbf{7}.
        \item Um atributo \texttt{int tamanho} (inicialmente 0) para rastrear a quantidade real de músicas.
    \end{itemize}
    \item Implemente o método \texttt{adicionarNoFim(String musica)} — insere no final da playlist. Caso a playlist esteja cheia, exiba uma mensagem de erro.
    \item Implemente o método \texttt{inserirNaPosicao(int posicao, String musica)} — insere uma música em uma posição específica, \textbf{deslocando (shift)} todos os elementos subsequentes uma posição para a direita. Valide se a posição é válida e se há espaço disponível.
    \item Implemente o método \texttt{removerDaPosicao(int posicao)} — remove a música na posição indicada, \textbf{deslocando (shift)} todos os elementos subsequentes uma posição para a esquerda para preencher o ``buraco''. Atribua \texttt{null} à última posição antes ocupada e decremente \texttt{tamanho}.
    \item Implemente o método \texttt{exibir()} que mostre apenas as músicas válidas (do índice 0 até \texttt{tamanho - 1}).
    \item Na classe \texttt{Main}, demonstre: adicione 5 músicas ao final, insira uma música na posição 2, remova a música na posição 0, e exiba a playlist após cada operação. Imprima \textbf{quantos deslocamentos (shifts)} foram necessários em cada inserção e remoção.
\end{itemize}

\textbf{Pergunta Reflexiva:} Ao inserir uma música na posição 0 de uma playlist que já contém $N$ músicas, todos os $N$ elementos precisam ser deslocados uma posição para a direita. Qual é a complexidade de tempo (Big-O) dessa operação no \textbf{pior caso}? Compare com a inserção no final. Por que essa diferença é uma limitação fundamental dos arrays primitivos?

\vspace{0.5cm}

\section*{Questão 6: Sistema de Senhas de Atendimento (Desafio)}
Um hospital utiliza um sistema de fila de atendimento baseado em senhas numéricas. Devido a uma restrição do sistema legado, a fila é implementada com um \textbf{array primitivo de inteiros} de tamanho fixo. O sistema precisa lidar com situações do mundo real como: pacientes que desistem (saem da fila), pacientes prioritários que precisam ser encaixados em posições específicas, e a possibilidade de a fila lotar. Além disso, o administrador quer uma opção de \textbf{redimensionamento manual} caso a capacidade se esgote.

\textbf{Requisitos:}
\begin{itemize}
    \item Crie uma classe \texttt{FilaDeSenhas} com:
    \begin{itemize}
        \item Um array \texttt{int[] senhas} com capacidade inicial de \textbf{5}.
        \item Um atributo \texttt{int quantidade} para controlar o número real de senhas na fila.
        \item Um atributo \texttt{int proximaSenha} (iniciando em 1) para gerar senhas sequenciais automaticamente.
    \end{itemize}
    \item Implemente \texttt{chamarProximaSenha()} — gera a próxima senha e a insere no \textbf{final} da fila. Se a fila estiver cheia, exiba ``Fila lotada!'' e \textbf{não} insira.
    \item Implemente \texttt{encaixarPrioritario(int posicao)} — gera uma nova senha e a insere na posição indicada, deslocando os demais. Valide se há espaço e se a posição é válida.
    \item Implemente \texttt{desistir(int posicao)} — remove a senha na posição informada, deslocando os elementos restantes para a esquerda.
    \item Implemente \texttt{redimensionar(int novaCapacidade)} — cria um \textbf{novo} array com a nova capacidade, copia todos os elementos existentes para ele e substitui o array interno. Trate o caso em que a nova capacidade é menor que a quantidade atual de elementos.
    \item Implemente \texttt{exibirFila()} para mostrar as senhas na ordem atual e a capacidade total.
    \item Na classe \texttt{Main}, simule o seguinte cenário:
    \begin{enumerate}
        \item Chame 5 senhas normais (lotando a fila).
        \item Tente chamar uma 6ª senha (deve mostrar ``Fila lotada!'').
        \item Redimensione a fila para capacidade 10.
        \item Encaixe um prioritário na posição 0.
        \item Faça a senha da posição 3 desistir.
        \item Exiba a fila final.
    \end{enumerate}
\end{itemize}

\textbf{Pergunta Reflexiva:} O método \texttt{redimensionar} cria um novo array e copia todos os elementos um a um. Sabendo que arrays ocupam \textbf{blocos contíguos de memória}, explique por que o Java (e a maioria das linguagens) não pode simplesmente ``esticar'' o array original na memória. O que acontece com o array antigo após a cópia? Relacione sua resposta com o conceito de \textit{Garbage Collection} da JVM.

\end{document}